\chapter{Machine Learning}

\begin{warningblock}[Work in Progress]
Chapter incomplete: content requires review.
\end{warningblock}

\section{Introduction}

Machine Learning (ML) provides a framework for tackling problems that are inherently difficult—if not impossible—to solve with conventional programming paradigms. Such problems often involve relationships between inputs and outputs that are highly nonlinear, high-dimensional, and not easily captured by analytical formulas. Typical examples include:

\begin{itemize}
    \item \textbf{Text classification}: Assign a label (e.g., sentiment analysis) to a text based on its content.
    \item \textbf{Fraud detection}: Estimate the probability that a credit card transaction is fraudulent.
    \item \textbf{Object recognition}: Identify objects within images, potentially from atypical viewpoints or under challenging conditions (e.g., varying illumination, occlusion, cluttered scenes).
\end{itemize}

In all these cases, the result hinges on a complex, nonlinear combination of many features or parameters, each individually contributing only incrementally to the final outcome.

\begin{observationblock}[Why ``Learning''?]
    \textit{Machine learning problems are, at their core, optimization problems: we seek parameters that minimize some discrepancy between prediction and reality. The learning aspect refers to the process by which these parameters are adjusted iteratively, based on patterns in observed data, especially in the absence of closed-form or mechanistic models.}
\end{observationblock}

\subsubsection*{A Canonical Machine Learning Workflow}
Suppose we are given a \textbf{training set} consisting of many input--output pairs $\{(x_k, y_k)\}$. The typical ML paradigm proceeds via:

\begin{enumerate}
    \item \textbf{Model specification}: Choose a parametric family of models $f(x;\Theta)$, where $\Theta$ denotes the set of adjustable parameters. This class should be flexible enough to capture the complexity of the mapping from $x$ to $y$.
    \item \textbf{Loss function}: Define a quantitative measure (loss function $\mathcal{L}(\Theta)$) that expresses how well the model’s predictions $f(x_k;\Theta)$ match the actual outputs $y_k$ in the data.
    \item \textbf{Optimization}: Adjust or ``train'' the parameters $\Theta$ to minimize the loss over the training data, i.e., to find
          $$
          \Theta^* = \arg\min_\Theta \mathcal{L}(\Theta)
          $$
\end{enumerate}

\begin{exampleblock}[Iterative Solutions]
    Unlike classical optimization, where an explicit, analytic solution may be available, the mapping learned in ML is rarely known in closed form. Instead, stochastic and iterative algorithms like gradient descent are used to progressively refine the parameters based on feedback from the data.
\end{exampleblock}

Machine learning thus describes a process in which models adaptively learn patterns from data through optimization, generalizing from past observations to accurately predict new, unseen data.

\subsubsection{Learning Tasks}

Machine learning encompasses a variety of learning paradigms suited to different problem types. The three main classes are:

\vspace{0.6em}
\begin{itemize}
    \item \textbf{Supervised learning}: Given a set of inputs and corresponding outputs (labels), the goal is to learn a mapping from inputs to outputs.
    \begin{itemize}
        \item \emph{Classification}: Predict discrete labels (e.g., cat vs.\ dog, spam vs.\ not-spam).
        \item \emph{Regression}: Predict continuous values (e.g., house prices, temperature).
    \end{itemize}
    % A typical supervised task is showing the model pairs $(x, y)$ and learning a function that maps $x \rightarrow y$ accurately.
    % \hfill
    % \includegraphics[height=2.5cm]{assets/supervised.png}
    \vspace{0.6em}

    \item \textbf{Unsupervised learning}: Only input data is provided—there are no explicit labels. The algorithm seeks to discover structure or patterns within the data.
    \begin{itemize}
        \item \emph{Clustering}: Group similar data points together.
        \item \emph{Component analysis}: Identify informative latent factors (e.g., PCA).
        \item \emph{Autoencoding}: Learn compact representations of the data.
    \end{itemize}
    % For example, points are grouped based on their distribution into natural clusters.
    % \hfill
    % \includegraphics[height=2.5cm]{assets/unsupervised.png}
    % \vspace{0.6em}

    \item \textbf{Reinforcement learning}: An agent interacts with an environment, taking actions and receiving feedback in the form of rewards. The goal is to learn a strategy (policy) that maximizes cumulative reward over time.
    \begin{itemize}
        \item \emph{Learning long-term gains}: Optimizing for actions that are beneficial in the long run, not just immediately.
        \item \emph{Planning}: Sequencing actions to achieve complex objectives.
    \end{itemize}
    % \hfill
    % \includegraphics[height=2.5cm]{assets/reinforcement.png}
\end{itemize}


\subsubsection{Classification vs. Regression}

\dots

\section{Neural Networks}

\subsection{Feedforward Neural Networks}

\subsubsection{Activation Functions}

\subsubsection{Training}

\subsubsection{Regularization}

\begin{itemize}
    \item Weight decay
    \item Dropout
\end{itemize}

\subsection{Convolutional Neural Networks}

\subsubsection{Convolutional Layers}

\begin{itemize}
    \item Padding
    \item Stride
    \item Pooling
    \item BatchNorm
\end{itemize}



\subsection{Normalizing Flows}

Normalizing Flows are a family of methods for constructing flexible learnable probability distributions,
often with neural networks, which allow us to surpass the limitations of simple parametric forms.
Key ideas:
● Learn a bijective mapping function f: R D
→RD
● Use change of variables formula to compute densities p(x) dx = p(z) dz → p(z) = p(x) |dx/dz|
Applications:
● Density Estimation: Exact likelihood modeling
● Bayesian Inference: Posterior sampling in simulation-based inference
● Image/Audio Generation: Realistic, high-resolution samples
● Anomaly Detection